\usepackage[utf8]{inputenc}
\usepackage[T1]{fontenc}
\usepackage[top=2.5cm, bottom=2.5cm, left=3cm, right=2.5cm]{geometry}
\usepackage{graphicx}
\usepackage{nopageno}
\usepackage{setspace}
\usepackage{amsmath}
\usepackage{titlesec}
\usepackage{array}
\usepackage{booktabs}
\usepackage{multirow}
\usepackage{float}
\usepackage{url}
\IfFileExists{newtxtext.sty}{%
    \usepackage{newtxtext,newtxmath}%
}{%
    \usepackage{mathptmx}%
}
\usepackage{multicol}
\usepackage[labelsep=period,skip=4pt]{caption}
\usepackage{indentfirst}
\usepackage{enumitem}
\usepackage{longtable}
\usepackage{tabularx}
\usepackage{siunitx}
\usepackage{pgfplots}
\usepackage{tikz}
\usepackage{xcolor}
\usepackage{listings}
\usepackage{hyperref}

\pgfplotsset{compat=1.16}
\usetikzlibrary{arrows.meta,positioning,calc}
\sisetup{
    output-decimal-marker={,},
    group-separator={.},
    group-digits=integer,
    range-phrase={--},
    range-units=single,
    detect-all,
    per-mode=symbol
}
\hypersetup{
    colorlinks=true,
    linkcolor=black,
    citecolor=black,
    urlcolor=black,
    hypertexnames=false,
    pdftitle={Laporan Audit Instrumentasi Group I - FIT310},
    pdfauthor={Group I}
}

% 1. Paragraph & spacing setting (Template O2)
\linespread{1.25}
\setlength{\parindent}{0.8cm}
\setstretch{1.25}
\setlength{\columnsep}{10pt}
\setlength{\parskip}{0pt}

\setlength{\textfloatsep}{6pt}
\setlength{\intextsep}{6pt}
\setlength{\floatsep}{6pt}

\titlespacing*{\section}{0pt}{8pt}{4pt}
\titlespacing*{\subsection}{0pt}{0.4em}{0.2em}
\titlespacing*{\subsubsection}{0pt}{0.3em}{0.15em}

\setlist[enumerate]{
    itemsep=2pt,
    parsep=0pt,
    topsep=0pt,
    partopsep=0pt
}
\setlist[itemize]{
    itemsep=2pt,
    parsep=0pt,
    topsep=0pt,
    partopsep=0pt
}

% 2. Figure and Table Setting
\renewcommand{\figurename}{Gambar}
\renewcommand{\tablename}{Tabel}

% 3. Pengaturan Section (Polosan, tanpa huruf/nomor) - Template O2
\renewcommand{\thesection}{}
\titleformat{\section}
{\normalfont\bfseries}{}{0pt}{}

% 4. Pengaturan Subsection (Format: a. Judul) - Template O2
\renewcommand{\thesubsection}{\alph{subsection}.}
\titleformat{\subsection}
{\normalfont\bfseries}{\thesubsection}{0.5em}{}

\renewcommand{\thesubsubsection}{}
\titleformat{\subsubsection}
{\normalfont\bfseries}{}{0pt}{}

% 5. Listing kode untuk Lampiran
\lstset{
    basicstyle=\ttfamily\footnotesize,
    breaklines=true,
    frame=single,
    columns=fullflexible,
    keepspaces=true,
    showstringspaces=false,
    xleftmargin=0.2cm
}

% 6. Makro metrologi
\newcommand{\uc}{u_{c}}
\newcommand{\uA}{u_{A}}
\newcommand{\uB}{u_{B}}
\newcommand{\dd}{\mathrm{d}}

% 7. Pustaka / Bibliografi (Template O2)
\usepackage{csquotes}
\IfFileExists{biblatex.sty}{%
    \usepackage[style=apa,sortcites=true,sorting=nyt,backend=biber]{biblatex}
    \DeclareLanguageMapping{american}{american-apa}
    \addbibresource{Pustaka.bib}
    \setlength{\bibitemsep}{0.0\baselineskip}
}{%
    \usepackage[numbers,sort&compress]{natbib}
    \bibliographystyle{apalike}
    \renewcommand{\bibsection}{}
    \setlength{\bibsep}{2pt}
    \newcommand{\printbibliography}[1][]{\footnotesize\bibliography{Pustaka}}
    \providecommand{\parencite}{\citep}
}
